\usepackage[margin=0.8in]{geometry}
\usepackage{amsmath,amssymb,amsthm}
\usepackage{enumitem}
\usepackage{longtable}
\usepackage{array}
\usepackage{xcolor}
\usepackage{hyperref}
\usepackage{fancyhdr}
\usepackage{titlesec}
\usepackage{tcolorbox}
\usepackage{multicol}
\usepackage{booktabs}
\hypersetup{colorlinks=true, linkcolor=blue!50!black, urlcolor=blue!50!black}
\pagestyle{fancy}
\fancyhf{}
\rhead{\thepage}
\lhead{Engineering Math to Pure Math}
\setlength{\parskip}{0.55em}
\setlength{\parindent}{0pt}
\setlist[itemize]{topsep=0.2em,itemsep=0.2em}
\setlist[enumerate]{topsep=0.2em,itemsep=0.25em}
\newtcolorbox{focusbox}{colback=blue!3!white,colframe=blue!50!black,title=Focus}
\newtcolorbox{intuitionbox}{colback=green!3!white,colframe=green!40!black,title=Intuition}
\newtcolorbox{formalbox}{colback=gray!4!white,colframe=gray!70!black,title=Formal}
\newtcolorbox{walkbox}{colback=orange!5!white,colframe=orange!60!black,title=Walking question}
\newcommand{\N}{\mathbb{N}}
\newcommand{\Z}{\mathbb{Z}}
\newcommand{\Q}{\mathbb{Q}}
\newcommand{\R}{\mathbb{R}}
\newcommand{\C}{\mathbb{C}}
\newcommand{\eps}{\varepsilon}
\newcommand{\pow}{\mathcal{P}}
\theoremstyle{definition}
\newtheorem{definition}{Definition}
\newtheorem{example}{Example}
\newtheorem{exercise}{Exercise}
\theoremstyle{plain}
\newtheorem{theorem}{Theorem}
\newtheorem{lemma}{Lemma}
\theoremstyle{remark}
\newtheorem{remark}{Remark}
